\documentclass[12pt]{beamer}
\usepackage{ceai-beamer}
\usepackage{ragged2e}
\AtBeginEnvironment{frame}{\justifying}

\title[Topic 11: Regression]{CE 488/5588: Applied AI\\Topic 11: Supervised ML for Regression}
\author{}
\date{}

\newcommand{\source}[1]{\vspace{0.4em}{\tiny\color{gray}Source: #1\par}}

\begin{document}

\begin{frame}
  \titlepage
\end{frame}

\begin{frame}{Lecture Roadmap}
  \justifying
  \tableofcontents
\end{frame}

\section{Regression as Supervised Prediction}

\subsection{Problem Setting and Linear Regression}

\begin{frame}{From Classification to Regression}
  \justifying
  In classification, the model predicts a category. In regression, the output itself is a real-valued quantity, so the numerical size of the prediction error matters directly.
  
  \begin{conceptbox}{Core Shift}
    The supervised-learning pipeline stays the same: features, labeled data, training, validation, and testing. What changes is the target type and how prediction quality is measured.
  \end{conceptbox}
  
  \source{\href{https://www.itl.nist.gov/div898/handbook/pmd/section1/pmd141.htm}{NIST e-Handbook: Linear Least Squares Regression}}
\end{frame}

\begin{frame}{Learning Goals for Topic 11}
  \justifying
  By the end of this topic, you should be able to:
  \begin{itemize}
    \item explain how regression fits within supervised learning;
    \item formulate linear regression in scalar and vectorized forms;
    \item derive the least-squares solution and explain regularization;
    \item evaluate models using MAE, MSE, RMSE, and $R^2$;
    \item describe nonlinear regression with MLPs and SVR; and
    \item cast time-series prediction as supervised regression with lag features.
  \end{itemize}
\end{frame}

\begin{frame}{Why Engineers Care About Regression}
  \justifying
  Many engineering tasks are naturally continuous rather than categorical.
  \begin{examplebox}{Typical Use Cases}
    \begin{itemize}
      \item estimating settlement or material strength,
      \item forecasting load, traffic speed, or reservoir level,
      \item approximating physical-system response, and
      \item predicting future sensor values in monitored systems.
    \end{itemize}
  \end{examplebox}
\end{frame}

\begin{frame}{Dataset and Prediction Mapping}
  \justifying
  Let the dataset be
  \[
    \mathcal{D} = \{(\mathbf{x}_i, y_i) \mid i=1,2,\ldots,N\},
  \]
  where $\mathbf{x}_i \in \mathbb{R}^{m}$ is the feature vector and $y_i \in \mathbb{R}$ is the target value.
  \[
    \hat{y} = f(\mathbf{x})
  \]
  is the learned mapping from inputs to a numerical output.
\end{frame}

\begin{frame}{Regression vs. Classification}
  \justifying
  \begin{columns}[T]
    \begin{column}{0.48\textwidth}
      \begin{conceptbox}{Classification}
        Small score changes may not alter the final predicted label.
      \end{conceptbox}
    \end{column}
    \begin{column}{0.48\textwidth}
      \begin{conceptbox}{Regression}
        The predicted value itself matters, so the error magnitude becomes central.
      \end{conceptbox}
    \end{column}
  \end{columns}
  \vspace{0.5em}
  In regression, model evaluation must reflect the operational cost of numerical error.
\end{frame}

\begin{frame}{Classical Linear Regression Model}
  \justifying
  We begin with the linear model
  \[
    f(\mathbf{x}) = a_1x_1 + a_2x_2 + \cdots + a_mx_m + a_0,
  \]
  where $a_1,\ldots,a_m$ are slope coefficients and $a_0$ is the intercept.
  
  \begin{conceptbox}{Why Start Here?}
    Linear regression gives the cleanest setting for defining residuals, objective functions, estimation, and model evaluation.
  \end{conceptbox}
  \source{\href{https://www.itl.nist.gov/div898/handbook/pmd/section1/pmd141.htm}{NIST e-Handbook}}
\end{frame}

\begin{frame}{Geometric Interpretation}
  \justifying
  The linear model defines a hyperplane in feature space.
  \begin{itemize}
    \item With one feature, the model is a line.
    \item With two features, it is a plane.
    \item In higher dimensions, it remains a flat affine surface.
  \end{itemize}
  Geometric thinking helps explain why linear regression is simple, interpretable, and limited when the true relationship is strongly nonlinear.
\end{frame}

\subsection{Residuals and Objective Functions}

\begin{frame}{Residuals Measure Prediction Error}
  \justifying
  For sample $i$, let
  \[
    \hat{y}_i = f(\mathbf{x}_i).
  \]
  The residual is
  \[
    e_i = y_i - \hat{y}_i = y_i - f(\mathbf{x}_i).
  \]
  A positive residual means the model under-predicted; a negative residual means it over-predicted.
\end{frame}

\begin{frame}{Residuals in One-Dimensional Linear Regression}
  \centering
  \begin{figure}
  \includegraphics[width=0.66\textwidth,height=0.44\textheight,keepaspectratio]{part03topic08fig01.png}
  \end{figure}
  \vspace{-0.6em}
  {\small Residuals are vertical distances between observed points and the fitted line.}
\end{frame}

\begin{frame}{From Residuals to an Objective Function}
  \justifying
  A model generates one residual per sample, so training requires an aggregate measure of error over all $N$ samples.
  
  \begin{conceptbox}{Estimation Problem}
    Choose parameters $\mathbf{a}$ that minimize a function of all residuals.
  \end{conceptbox}
  Different objectives emphasize errors differently and therefore produce different fitted models.
\end{frame}

\begin{frame}{Sum of Absolute Errors (SAE)}
  \justifying
  One common choice is the sum of absolute errors:
  \[
    E_1(\mathbf{a}) = \sum_{i=1}^{N} \left|y_i - f(\mathbf{x}_i)\right|.
  \]
  \begin{conceptbox}{Interpretation}
    SAE grows linearly with error magnitude, so it is less sensitive to large outliers than squared loss.
  \end{conceptbox}
\end{frame}

\begin{frame}{Sum of Squared Errors (SSE)}
  \justifying
  Another standard choice is the sum of squared errors:
  \[
    E_2(\mathbf{a}) = \sum_{i=1}^{N} \left(y_i - f(\mathbf{x}_i)\right)^2.
  \]
  \begin{conceptbox}{Interpretation}
    SSE penalizes large errors quadratically, which makes the model much more sensitive to large misses.
  \end{conceptbox}
\end{frame}

\begin{frame}{SAE vs. SSE}
  \justifying
  \begin{columns}[T]
    \begin{column}{0.48\textwidth}
      \begin{examplebox}{SAE}
        \begin{itemize}
          \item robust to outliers,
          \item penalty grows linearly,
          \item objective is convex but not differentiable at zero.
        \end{itemize}
      \end{examplebox}
    \end{column}
    \begin{column}{0.48\textwidth}
      \begin{warningbox}{SSE}
        \begin{itemize}
          \item heavily penalizes large errors,
          \item smooth objective,
          \item leads to a closed-form solution for linear models.
        \end{itemize}
      \end{warningbox}
    \end{column}
  \end{columns}
\end{frame}

\begin{frame}{Why Least Squares Is Often Preferred}
  \justifying
  For linear models, the squared-error objective is especially attractive because it is smooth and differentiable.
  \[
    \hat{\mathbf{a}} = \underset{\mathbf{a}}{\operatorname{argmin}}\; E(\mathbf{a})
  \]
  becomes analytically tractable when $E(\mathbf{a}) = \lVert \mathbf{y} - \mathbf{X}\mathbf{a} \rVert_2^2$.
  
  \source{\href{https://www.itl.nist.gov/div898/handbook/pmd/section4/pmd431.htm}{NIST e-Handbook: Least Squares}; \href{https://cs229.stanford.edu/main_notes.pdf}{Stanford CS229 Notes}}
\end{frame}

\subsection{Vectorized Least-Squares Formulation}

\begin{frame}{Augmented Feature Vector}
  \justifying
  Introduce the augmented vector
  \[
    \tilde{\mathbf{x}} = [x_1, x_2, \ldots, x_m, 1]^T,
  \]
  so the model becomes
  \[
    f(\mathbf{x}) = \tilde{\mathbf{x}}^T\mathbf{a}.
  \]
  This moves the intercept into the parameter vector and simplifies matrix notation.
\end{frame}

\begin{frame}{Design Matrix and Target Vector}
  \justifying
  Stack all augmented feature vectors into the design matrix
  \[
    \mathbf{X} =
    \begin{bmatrix}
      \tilde{\mathbf{x}}_1^T \\
      \tilde{\mathbf{x}}_2^T \\
      \vdots \\
      \tilde{\mathbf{x}}_N^T
    \end{bmatrix}
    \in \mathbb{R}^{N \times (m+1)},
  \]
  and the targets into
  \[
    \mathbf{y} = [y_1, y_2, \ldots, y_N]^T.
  \]
\end{frame}

\begin{frame}{Vectorized Residual Objectives}
  \justifying
  The prediction vector is
  \[
    \hat{\mathbf{y}} = \mathbf{X}\mathbf{a}.
  \]
  Therefore,
  \[
    E_1(\mathbf{a}) = \lVert \mathbf{y} - \mathbf{X}\mathbf{a} \rVert_1,
    \qquad
    E_2(\mathbf{a}) = \lVert \mathbf{y} - \mathbf{X}\mathbf{a} \rVert_2^2.
  \]
  Vector notation reveals regression as optimization over a norm of the residual vector.
\end{frame}

\begin{frame}{Least-Squares Objective}
  \justifying
  The least-squares problem is
  \[
    \min_{\mathbf{a}}\; E(\mathbf{a}) = \lVert \mathbf{y} - \mathbf{X}\mathbf{a} \rVert_2^2.
  \]
  Expanding the objective gives
  \[
    E(\mathbf{a}) = (\mathbf{y} - \mathbf{X}\mathbf{a})^T(\mathbf{y} - \mathbf{X}\mathbf{a}).
  \]
  This quadratic form is the basis of the normal-equation derivation.
\end{frame}

\begin{frame}{Gradient and Normal Equations}
  \justifying
  Taking the gradient with respect to $\mathbf{a}$ yields
  \[
    \nabla E(\mathbf{a}) = -2\mathbf{X}^T(\mathbf{y} - \mathbf{X}\mathbf{a}).
  \]
  Setting the gradient to zero gives the normal equations:
  \[
    \mathbf{X}^T\mathbf{X}\mathbf{a} = \mathbf{X}^T\mathbf{y}.
  \]
  \source{\href{https://cs229.stanford.edu/main_notes.pdf}{Stanford CS229 Notes}; \href{https://www.itl.nist.gov/div898/handbook/pmd/section4/pmd431.htm}{NIST Least Squares}}
\end{frame}

\begin{frame}{Closed-Form Least-Squares Solution}
  \justifying
  If $\mathbf{X}^T\mathbf{X}$ is invertible, then
  \[
    \hat{\mathbf{a}} = (\mathbf{X}^T\mathbf{X})^{-1}\mathbf{X}^T\mathbf{y}.
  \]
  \begin{itemize}
    \item $\mathbf{X}$ is $N \times (m+1)$,
    \item $\mathbf{X}^T\mathbf{X}$ is $(m+1) \times (m+1)$,
    \item $\hat{\mathbf{a}}$ is a vector in $\mathbb{R}^{m+1}$.
  \end{itemize}
\end{frame}

\begin{frame}{When Ordinary Least Squares Becomes Unstable}
  \justifying
  Problems arise when $\mathbf{X}^T\mathbf{X}$ is singular or ill-conditioned.
  \begin{warningbox}{Common Causes}
    \begin{itemize}
      \item strongly correlated features,
      \item too many features relative to data,
      \item noisy data and over-flexible models.
    \end{itemize}
  \end{warningbox}
  These issues motivate regularization.
\end{frame}

\subsection{Regularized Linear Regression and Evaluation}

\begin{frame}{Regularization Adds a Complexity Penalty}
  \justifying
  Regularization augments the data-fit objective with a penalty on coefficient size.
  \begin{conceptbox}{Main Idea}
    Fit the training data well, but discourage unnecessarily large coefficients that may amplify noise or numerical instability.
  \end{conceptbox}
  It improves stability and can improve generalization.
\end{frame}

\begin{frame}{Ridge vs. Lasso Regression}
  \justifying
  \[
    \text{Ridge: } \min_{\mathbf{a}} \; \lVert \mathbf{y} - \mathbf{X}\mathbf{a} \rVert_2^2 + \lambda \lVert \mathbf{a} \rVert_2^2
  \]
  \[
    \text{Lasso: } \min_{\mathbf{a}} \; \lVert \mathbf{y} - \mathbf{X}\mathbf{a} \rVert_2^2 + \lambda \lVert \mathbf{a} \rVert_1
  \]
  \begin{itemize}
    \item Ridge shrinks coefficients smoothly.
    \item Lasso can drive some coefficients exactly to zero.
  \end{itemize}
  \source{\href{https://scikit-learn.org/stable/modules/generated/sklearn.linear_model.Ridge.html}{scikit-learn Ridge}; \href{https://scikit-learn.org/stable/modules/generated/sklearn.linear_model.Lasso.html}{scikit-learn Lasso}}
\end{frame}

\begin{frame}{The Role of $\lambda$}
  \justifying
  The regularization parameter $\lambda$ controls the tradeoff between data fit and model complexity.
  \begin{columns}[T]
    \begin{column}{0.48\textwidth}
      \begin{examplebox}{Large $\lambda$}
        More shrinkage, lower variance, higher bias.
      \end{examplebox}
    \end{column}
    \begin{column}{0.48\textwidth}
      \begin{examplebox}{Small $\lambda$}
        Better training fit, higher variance, greater risk of overfitting.
      \end{examplebox}
    \end{column}
  \end{columns}
\end{frame}

\begin{frame}{How Should We Evaluate a Regression Model?}
  \justifying
  Training error alone is not enough. We care about performance on validation or test data.
  \begin{conceptbox}{Evaluation Principle}
    A useful metric should match the application: its units, its tolerance for outliers, and the practical cost of prediction error.
  \end{conceptbox}
\end{frame}

\begin{frame}{Common Regression Metrics}
  \justifying
  \[
    \mathrm{MAE} = \frac{1}{N} \sum_{i=1}^{N} |y_i - \hat{y}_i|,
    \qquad
    \mathrm{MSE} = \frac{1}{N} \sum_{i=1}^{N} (y_i - \hat{y}_i)^2
  \]
  \[
    \mathrm{RMSE} = \sqrt{\mathrm{MSE}},
    \qquad
    R^2 = 1 - \frac{\sum_{i=1}^{N}(y_i-\hat{y}_i)^2}{\sum_{i=1}^{N}(y_i-\bar{y})^2}
  \]
  \source{\href{https://scikit-learn.org/stable/modules/generated/sklearn.metrics.mean_absolute_error.html}{scikit-learn MAE}; \href{https://scikit-learn.org/stable/modules/generated/sklearn.metrics.mean_squared_error.html}{MSE}; \href{https://scikit-learn.org/stable/modules/generated/sklearn.metrics.root_mean_squared_error.html}{RMSE}; \href{https://scikit-learn.org/stable/modules/generated/sklearn.metrics.r2_score.html}{$R^2$}}
\end{frame}

\begin{frame}{Choosing a Metric in Engineering Context}
  \justifying
  \begin{itemize}
    \item MAE is easy to interpret and robust to outliers.
    \item MSE and RMSE emphasize large errors.
    \item $R^2$ summarizes explained variation relative to predicting the mean.
  \end{itemize}
  \begin{warningbox}{No Universal Winner}
    The right metric depends on how costly large misses are and whether interpretability in physical units matters.
  \end{warningbox}
  \source{\href{https://online.stat.psu.edu/stat501/lesson/1/1.5}{PSU STAT 501 on $R^2$}}
\end{frame}

\section{Neural-Network Regression with Multi-Layer Perceptrons}

\subsection{Basic Idea and Training}

\begin{frame}{Why Move Beyond Linear Regression?}
  \justifying
  Linear regression is simple and interpretable, but many engineering systems are nonlinear.
  
  \begin{conceptbox}{MLP Motivation}
    By composing affine transformations with nonlinear activation functions, a Multi-Layer Perceptron can represent curved trends and variable interactions that a single hyperplane cannot.
  \end{conceptbox}
\end{frame}

\begin{frame}{Single-Hidden-Layer MLP Regressor}
  \justifying
  For one hidden layer,
  \[
    \hat{y} = \mathbf{w}_2^T\,\sigma(\mathbf{W}_1\mathbf{x}+\mathbf{b}_1) + b_2.
  \]
  The hidden layer introduces nonlinearity through $\sigma(\cdot)$, while the output layer is typically linear for regression.
  \source{\href{https://scikit-learn.org/stable/modules/neural_networks_supervised.html}{scikit-learn Neural Networks}; \href{https://scikit-learn.org/stable/modules/generated/sklearn.neural_network.MLPRegressor.html}{MLPRegressor}}
\end{frame}

\begin{frame}{MLPs Can Learn Nonlinear Trends}
  \centering
  \begin{figure}
    \includegraphics[width=0.72\textwidth,height=0.44\textheight,keepaspectratio]{sin-mlp.pdf}
  \end{figure}
  \vspace{-0.6em}
  {\small An MLP regressor can approximate curved patterns such as sinusoidal behavior.}
\end{frame}

\begin{frame}{Training Objective and Optimization}
  \justifying
  MLP regression is commonly trained by minimizing mean squared error:
  \[
    \min_{\boldsymbol{\theta}}\; \frac{1}{N}\sum_{i=1}^{N}\left(y_i - f_{\boldsymbol{\theta}}(\mathbf{x}_i)\right)^2.
  \]
  In practice, optimization uses backpropagation with algorithms such as stochastic gradient descent or Adam.
\end{frame}

\begin{frame}{Generalization Controls for MLPs}
  \justifying
  \begin{warningbox}{Overfitting Risk}
    Because MLPs are flexible, they can fit noise when data are limited or training is poorly controlled.
  \end{warningbox}
  \begin{itemize}
    \item use a validation set and early stopping,
    \item apply L2 regularization,
    \item limit network size when data are scarce,
    \item judge success on held-out data, not training loss alone.
  \end{itemize}
\end{frame}

\begin{frame}{MLP Regression: Strengths and Tradeoffs}
  \justifying
  \begin{columns}[T]
    \begin{column}{0.48\textwidth}
      \begin{examplebox}{Strengths}
        Flexible nonlinear approximation, interaction modeling, broad applicability.
      \end{examplebox}
    \end{column}
    \begin{column}{0.48\textwidth}
      \begin{warningbox}{Tradeoffs}
        Harder training, more tuning, weaker interpretability than ordinary linear regression.
      \end{warningbox}
    \end{column}
  \end{columns}
\end{frame}

\section{Support Vector Regression (SVR)}

\subsection{Linear SVR, Optimization, and Kernels}

\begin{frame}{Support Vector Regression Intuition}
  \justifying
  SVR adapts margin-based ideas from support vector machines to continuous targets.
  \[
    f(\mathbf{x}) = \mathbf{w}^T\mathbf{x} + b
  \]
  for the linear case.
  
  \begin{conceptbox}{Key Goal}
    Find a function that is accurate while remaining as flat as possible.
  \end{conceptbox}
  \source{\href{https://scikit-learn.org/stable/modules/generated/sklearn.svm.SVR.html}{scikit-learn SVR}}
\end{frame}

\begin{frame}{The $\epsilon$-Insensitive Loss}
  \justifying
  SVR ignores small residuals inside an acceptable tube:
  \[
    L_{\epsilon}(y,f(\mathbf{x})) = \max\left(0,\,|y-f(\mathbf{x})|-\epsilon\right).
  \]
  \begin{conceptbox}{Interpretation}
    Errors smaller than $\epsilon$ are treated as acceptable; only larger deviations incur loss.
  \end{conceptbox}
  \source{\href{https://scikit-learn.org/stable/modules/svm.html}{scikit-learn SVM User Guide}}
\end{frame}

\begin{frame}{Primal Optimization Problem in SVR}
  \justifying
  The standard primal SVR problem is
  \[
    \min_{\mathbf{w},b,\xi_i,\xi_i^*}\; \frac{1}{2}\lVert\mathbf{w}\rVert_2^2 + C\sum_{i=1}^{N}(\xi_i+\xi_i^*),
  \]
  where $C$ balances flatness and tolerance for errors outside the $\epsilon$-tube.
\end{frame}

\begin{frame}{Slack Variables and Constraints}
  \justifying
  The optimization is subject to
  \[
    y_i - (\mathbf{w}^T\mathbf{x}_i + b) \leq \epsilon + \xi_i,
  \]
  \[
    (\mathbf{w}^T\mathbf{x}_i + b) - y_i \leq \epsilon + \xi_i^*,
    \qquad \xi_i,\xi_i^* \ge 0.
  \]
  Slack variables let some points fall outside the tube while still keeping the optimization convex.
\end{frame}

\begin{frame}{Kernel SVR for Nonlinear Regression}
  \justifying
  Like SVM classification, SVR can be kernelized:
  \[
    k(\mathbf{x}_i,\mathbf{x}_j) = \langle \phi(\mathbf{x}_i), \phi(\mathbf{x}_j) \rangle,
  \]
  leading to
  \[
    f(\mathbf{x}) = \sum_{i=1}^{N}(\alpha_i-\alpha_i^*)k(\mathbf{x}_i,\mathbf{x}) + b.
  \]
  Only support vectors contribute nonzero coefficients.
  \source{\href{https://scikit-learn.org/stable/modules/svm.html}{scikit-learn SVM Guide}}
\end{frame}

\begin{frame}{Kernel SVR Example}
  \centering
  \begin{figure}
    \includegraphics[width=0.72\textwidth,height=0.44\textheight,keepaspectratio]{cos-svr.pdf}
  \end{figure}
  \vspace{-0.6em}
  {\small Kernel SVR can represent nonlinear trends while retaining margin-based regularization.}
\end{frame}

\begin{frame}{When SVR Is a Good Choice}
  \justifying
  SVR is often effective when the dataset is moderate in size and strong regularization is desirable.
  \begin{itemize}
    \item $C$ controls the fit--flatness tradeoff,
    \item $\epsilon$ sets the tolerance band,
    \item kernel parameters shape nonlinear flexibility.
  \end{itemize}
\end{frame}

\section{Time-Series Prediction as Supervised Regression}

\subsection{Lag Features and Engineering Examples}

\begin{frame}{Forecasting Can Be Framed as Regression}
  \justifying
  Many engineering datasets are time-indexed: load, traffic speed, reservoir level, and vibration measurements all evolve over time.
  
  \begin{conceptbox}{Key Viewpoint}
    Many forecasting problems become ordinary supervised regression once past observations are reorganized into input features and future values are used as targets.
  \end{conceptbox}
\end{frame}

\begin{frame}{Lag Features and Forecast Targets}
  \justifying
  If $s_t$ is a scalar signal observed at time $t$, then we can build
  \[
    \mathbf{x}_t = [s_{t-1}, s_{t-2}, \ldots, s_{t-p}]^T,
  \]
  and define the target as
  \[
    y_t = s_t \qquad \text{or} \qquad y_t = s_{t+h}.
  \]
  Additional exogenous variables such as weather, occupancy, or control inputs can be appended to the feature vector.
  \source{\href{https://otexts.com/fpp3/AR.html}{Forecasting: Principles and Practice}; \href{https://scikit-learn.org/stable/auto_examples/applications/plot_time_series_lagged_features.html}{scikit-learn lagged-feature example}}
\end{frame}

\begin{frame}{Engineering Forecasting Example}
  \centering
  \begin{figure}
    \includegraphics[width=0.8\textwidth,height=0.42\textheight,keepaspectratio]{topic11_load_forecast.png}
  \end{figure}
  \vspace{-0.6em}
  {\small Load forecasting can be modeled with lagged demand, temperature, and calendar features.}
\end{frame}

\begin{frame}{A Time-Series Caution: Respect Temporal Order}
  \justifying
  The supervised-regression viewpoint does not remove time dependence.
  \begin{warningbox}{Validation Rule}
    Random train/test shuffles can be overly optimistic for forecasting tasks. Validation splits should respect chronology.
  \end{warningbox}
  We still learn a mapping from inputs to a numerical target, but feature construction and data splitting must respect time.
  \source{\href{https://scikit-learn.org/stable/auto_examples/applications/plot_time_series_lagged_features.html}{scikit-learn lagged-feature example}}
\end{frame}

\section{Conclusion Remarks}

\begin{frame}{Key Takeaways from Topic 11}
  \justifying
  \begin{itemize}
    \item Regression predicts continuous quantities rather than class labels.
    \item Linear regression introduces residuals, objective functions, and least-squares estimation.
    \item Regularization improves stability and controls overfitting.
    \item MLPs and SVR provide nonlinear alternatives with different tradeoffs.
    \item Forecasting problems can often be converted into supervised regression with lag features.
  \end{itemize}
\end{frame}

\begin{frame}{Choosing the Right Regressor}
  \justifying
  Model choice should reflect the structure of the data and the cost of error.
  \begin{columns}[T]
    \begin{column}{0.32\textwidth}
      \begin{examplebox}{Linear Models}
        Simple, interpretable, often strong baselines.
      \end{examplebox}
    \end{column}
    \begin{column}{0.32\textwidth}
      \begin{examplebox}{MLPs}
        Flexible and expressive for nonlinear relationships.
      \end{examplebox}
    \end{column}
    \begin{column}{0.32\textwidth}
      \begin{examplebox}{SVR}
        Strong regularization and geometric intuition.
      \end{examplebox}
    \end{column}
  \end{columns}
\end{frame}

\end{document}
